\documentclass[a4paper,11pt]{article}

\usepackage[utf8]{inputenc}
\usepackage[T1]{fontenc}
\usepackage{geometry}
\usepackage{longtable}
\usepackage{booktabs}
\usepackage{xcolor}
\usepackage{hyperref}
\usepackage{listings}
\usepackage{enumitem}

\geometry{margin=2.5cm}

\lstset{
	basicstyle=\ttfamily\small,
	breaklines=true,
	frame=single,
	backgroundcolor=\color{gray!5}
}

\title{Interface DISPO \texorpdfstring{$\leftrightarrow$}{↔} Backend}
\author{}
\date{}

\begin{document}
	
	\maketitle
	
	\section{Purpose}
	
	This document defines the REST API contract between the external DISPO system and the Rückbestätigungs-Backend.
	
	The DISPO system remains the leading system for order data and delivery planning.  
	The backend does not create business orders independently. It receives confirmation-relevant order data, stores a local snapshot, creates confirmation requests, sends customer notifications and manages confirmation status.
	
	\section{System Boundary}
	
	\subsection{DISPO System Responsibilities}
	
	\begin{itemize}
		\item managing existing orders
		\item delivery planning
		\item assigning delivery dates and delivery windows
		\item dispatcher workflow
		\item manual start of confirmation process
	\end{itemize}
	
	\subsection{Backend Responsibilities}
	
	\begin{itemize}
		\item receiving confirmation requests
		\item validating incoming order data
		\item storing local order snapshots
		\item creating confirmation requests
		\item generating token-based action URL
		\item sending customer notification e-mails
		\item managing confirmation status
		\item storing customer responses
		\item handling expiration after 24 hours
	\end{itemize}
	
	\section{Process Start}
	
	The dispatcher manually starts the process inside DISPO:
	
	\begin{lstlisting}
		Ruckbestatigung senden
	\end{lstlisting}
	
	DISPO then calls the backend via REST API.
	
	\section{REST Endpoint}
	
	\subsection{Create Confirmation Request}
	
	\begin{lstlisting}
		POST /api/dispo/confirmation-requests
	\end{lstlisting}
	
	This endpoint creates a confirmation request and triggers synchronous e-mail sending.
	
	\section{Request DTO}
	
	\subsection{Request JSON}
	
	\begin{lstlisting}
		{
			"externalOrderId": "A-1024",
			"customerName": "Max Muller",
			"customerEmail": "daniel@example.com",
			"deliveryAddress": "Beispielstrase 12, 97070 Wurzburg",
			"product": "Heizol",
			"quantityLiters": 3000,
			"deliveryDate": "2026-06-12",
			"deliveryWindowStart": "10:00",
			"deliveryWindowEnd": "11:00"
		}
	\end{lstlisting}
	
	\subsection{Field Description}
	
	\begin{longtable}{p{4cm}p{2cm}p{2cm}p{6cm}}
		\toprule
		Field & Type & Required & Description \\
		\midrule
		externalOrderId & string & yes & Order identifier from DISPO \\
		customerName & string & yes & Customer name \\
		customerEmail & string & yes & E-mail address for confirmation \\
		deliveryAddress & string & yes & Full delivery address \\
		product & string & yes & Product to deliver \\
		quantityLiters & integer & yes & Ordered quantity in liters \\
		deliveryDate & date & yes & ISO date YYYY-MM-DD \\
		deliveryWindowStart & time & yes & HH:mm \\
		deliveryWindowEnd & time & yes & HH:mm \\
		\bottomrule
	\end{longtable}
	
	\subsection{Fields intentionally excluded in v1}
	
	\begin{itemize}
		\item externalTourId
		\item customerNumber
		\item pricePer100Liters
		\item dispatcherId
		\item dispatcherName
	\end{itemize}
	
\section{Success Response}

\subsection{201 Created}

Returned only if:

\begin{itemize}
	\item request validation succeeded
	\item confirmation request was created internally
	\item token URLs were generated internally
	\item e-mail was sent successfully
	\item status was set to SENT
\end{itemize}

\begin{lstlisting}
	{
		"externalOrderId": "A-1024",
		"confirmationStatus": "SENT"
	}
\end{lstlisting}

\subsection{Response Fields}

\begin{longtable}{p{4cm}p{2cm}p{8cm}}
	\toprule
	Field & Type & Description \\
	\midrule
	externalOrderId & string & DISPO order identifier used for correlation \\
	confirmationStatus & string & Current confirmation status, normally SENT \\
	\bottomrule
\end{longtable}

\subsection{Internal Fields not exposed in v1}

The backend internally creates and stores confirmation requests. However, the internal confirmation request id is not exposed to DISPO in API v1.

Internally generated but not returned:

\begin{itemize}
	\item confirmationRequestId
	\item confirmActionUrl
	\item rejectActionUrl
	\item sentAt
	\item expiresAt
\end{itemize}
	
	\section{Confirmation Status Values}
	
	\begin{longtable}{p{3cm}p{3cm}p{8cm}}
		\toprule
		API Value & German Label & Meaning \\
		\midrule
		SENT & versendet & Waiting for customer response \\
		CONFIRMED & bestätigt & Customer accepted delivery window \\
		REJECTED & abgelehnt & Customer rejected delivery window \\
		NO\_RESPONSE & keine Rückmeldung & No valid response within 24h \\
		\bottomrule
	\end{longtable}
	
\subsection{Design Decision: No confirmationRequestId in DISPO API v1}

The DISPO-facing API communicates on order level.

DISPO identifies the order through externalOrderId. Individual confirmation requests, token handling, active request tracking, old request invalidation and expiration handling are managed internally by the backend.

Therefore, confirmationRequestId is not exposed to DISPO in API v1.

This avoids coupling DISPO to the internal data model of the confirmation backend.

The backend is responsible for ensuring that only the latest active confirmation request can update the current confirmation status of an order.
	
\section{Validation Rules}

\begin{longtable}{p{2cm}p{4cm}p{8cm}}
	\toprule
	Rule ID & Field & Rule \\
	\midrule
	VAL-001 & externalOrderId & must not be blank \\
	VAL-002 & customerName & must not be blank \\
	VAL-003 & customerEmail & must not be blank and must be a valid e-mail address \\
	VAL-004 & deliveryAddress & must not be blank \\
	VAL-005 & product & must not be blank \\
	VAL-006 & quantityLiters & must be greater than 0 \\
	VAL-007 & deliveryDate & must be a valid date in ISO format YYYY-MM-DD \\
	VAL-008 & deliveryWindowStart & must be a valid time in format HH:mm \\
	VAL-009 & deliveryWindowEnd & must be a valid time in format HH:mm \\
	VAL-010 & deliveryWindow & deliveryWindowStart must be before deliveryWindowEnd \\
	\bottomrule
\end{longtable}
	
\section{Error Responses}

All error responses use the same structure.

\begin{lstlisting}
	{
		"code": "VALIDATION_ERROR",
		"message": "Human-readable error message",
		"status": 400,
		"path": "/api/dispo/confirmation-requests",
		"timestamp": "2026-06-12T08:15:00Z"
	}
\end{lstlisting}

\subsection{Fields}

\begin{longtable}{p{3cm}p{2cm}p{8cm}}
	\toprule
	Field & Type & Description \\
	\midrule
	code & string & Stable machine-readable error code \\
	message & string & Human-readable error description \\
	status & integer & HTTP status code \\
	path & string & Request path where the error occurred \\
	timestamp & datetime & Time when the error occurred \\
	\bottomrule
\end{longtable}

\subsection{400 Validation Error}

Returned when the request contains invalid or missing data.

\begin{lstlisting}
	400 Bad Request
\end{lstlisting}

\begin{lstlisting}
	{
		"code": "VALIDATION_ERROR",
		"message": "Customer e-mail is required for digital confirmation.",
		"status": 400,
		"path": "/api/dispo/confirmation-requests",
		"timestamp": "2026-06-12T08:15:00Z"
	}
\end{lstlisting}

\subsection{502 E-Mail Sending Failed}

Returned when the confirmation e-mail could not be sent.

\begin{lstlisting}
	502 Bad Gateway
\end{lstlisting}

\begin{lstlisting}
	{
		"code": "EMAIL_SENDING_FAILED",
		"message": "The confirmation e-mail could not be sent.",
		"status": 502,
		"path": "/api/dispo/confirmation-requests",
		"timestamp": "2026-06-12T08:15:00Z"
	}
\end{lstlisting}

\subsection{500 Internal Error}

Returned when an unexpected technical error occurs.

\begin{lstlisting}
	500 Internal Server Error
\end{lstlisting}

\begin{lstlisting}
	{
		"code": "INTERNAL_ERROR",
		"message": "An unexpected technical error occurred.",
		"status": 500,
		"path": "/api/dispo/confirmation-requests",
		"timestamp": "2026-06-12T08:15:00Z"
	}
\end{lstlisting}

\section{Business Rules}

\begin{longtable}{p{2cm}p{11cm}}
	\toprule
	Rule ID & Rule \\
	\midrule
	BR-001 & The order already exists in DISPO before the confirmation process starts. \\
	BR-002 & DISPO is the leading system for order data and delivery planning. \\
	BR-003 & The backend stores a local snapshot only when DISPO starts a confirmation request. \\
	BR-004 & The confirmation process is started manually by the dispatcher in DISPO. \\
	BR-005 & DISPO calls the backend directly via REST API. \\
	BR-006 & A confirmation request always refers to exactly one order and one concrete delivery window. \\
	BR-007 & Tour data is not part of API v1. \\
	BR-008 & The first successful confirmation status returned by the backend is SENT. \\
	BR-009 & SENT is set only after successful technical e-mail sending. \\
	BR-010 & The 24-hour response timer starts only after successful technical e-mail sending. \\
	BR-011 & If e-mail sending fails, the API returns an error and no successful confirmation status is created. \\
	BR-012 & The e-mail contains one token-based customer link. \\
	BR-013 & Action links open a minimal customer confirmation page. \\
	BR-014 & Opening an action link does not directly change the confirmation status. \\
	BR-015 & The confirmation status is changed only after a final submit on the customer page. \\
	BR-016 & Only one valid customer response is allowed per confirmation request. \\
	BR-017 & Only the latest active request can update the current confirmation status. \\
	BR-018 & Old, inactive or replaced links must not update the current order status. \\
	BR-019 & Customer comments are optional. \\
	BR-020 & Customer comments are stored but not automatically interpreted. \\
	BR-021 & The system does not perform automatic tour planning, route optimization or replanning. \\
	BR-022 & If the customer rejects or does not respond, the dispatcher decides manually how to proceed. \\
	BR-023 & When the status changes to CONFIRMED, REJECTED or NO\_RESPONSE, the backend sends a callback to DISPO. \\
	BR-024 & If DISPO sends the same request for the same externalOrderId with unchanged confirmation-relevant data, the backend treats it as duplicate and does not send a second e-mail. \\
	BR-025 & If DISPO sends a request for the same externalOrderId with changed confirmation-relevant data, the backend invalidates previous active requests and creates a new confirmation request. \\
	\bottomrule
\end{longtable}
	
	\section{Customer Confirmation Link Concept}
	
	The e-mail contains one token-based customer link:
	
	\begin{lstlisting}
		/confirmation/{token}
	\end{lstlisting}
	
	Opening this link does not change the confirmation status.
	
	The link opens a minimal customer confirmation page. On this page, the customer can choose between confirming or rejecting the proposed delivery window.
	
	\begin{lstlisting}
		GET /confirmation/{token} -> open customer page
		POST confirm action      -> update status to CONFIRMED
		POST reject action       -> update status to REJECTED
	\end{lstlisting}
	
	Reason: GET requests must not change state. E-mail scanners or preview systems may open links automatically.
	
	
	\section{Duplicate Handling}
	
	DISPO may send more than one request for the same externalOrderId.
	
	The backend distinguishes between duplicate requests with unchanged confirmation data and new requests with changed confirmation-relevant data.
	
	\subsection{Duplicate with unchanged data}
	
	If DISPO sends the same request again and the confirmation-relevant data has not changed, the backend treats the request as duplicate.
	
	Confirmation-relevant data includes:
	
	\begin{itemize}
		\item customerEmail
		\item deliveryAddress
		\item product
		\item quantityLiters
		\item deliveryDate
		\item deliveryWindowStart
		\item deliveryWindowEnd
	\end{itemize}
	
	In this case:
	
	\begin{itemize}
		\item no new confirmation request is created
		\item no second e-mail is sent
		\item the existing active confirmation request is returned
	\end{itemize}
	
	\begin{lstlisting}
		200 OK
	\end{lstlisting}
	
	\begin{lstlisting}
		{
			"externalOrderId": "A-1024",
			"confirmationStatus": "SENT"
		}
	\end{lstlisting}
	
	\subsection{Changed confirmation-relevant data}
	
	If DISPO sends a request for the same externalOrderId but confirmation-relevant data has changed, the backend treats it as a new confirmation request.
	
	In this case:
	
	\begin{itemize}
		\item previous active requests are invalidated internally
		\item a new confirmation request is created
		\item a new e-mail is sent
		\item the new confirmation request is returned
	\end{itemize}
	
	\begin{lstlisting}
		201 Created
	\end{lstlisting}
	
	\begin{lstlisting}
		{
			"externalOrderId": "A-1024",
			"confirmationStatus": "SENT"
		}
	\end{lstlisting}
	
	The API v1 does not return 409 Conflict for duplicate active requests.
	
	\section{Sequence Overview}
	
	\subsection{Successful Flow}
	
	\begin{lstlisting}
		DISPO -> Backend -> Notification Service -> Backend -> DISPO
	\end{lstlisting}
	
	\subsection{Failed E-Mail Flow}
	
	\begin{lstlisting}
		DISPO -> Backend -> Notification Service (fail) -> Backend -> 502
	\end{lstlisting}
	
	\section{Assumptions}
	
	\begin{itemize}
		\item DISPO can call backend via REST
		\item e-mail is MVP channel
		\item sending is synchronous
		\item expiration is 24h
	\end{itemize}
	
	\section{Open Questions}
	
	\begin{itemize}
		\item Which exact fields does real DISPO provide?
		\item Should callback retries be implemented?
		\item Should reminder logic be added?
	\end{itemize}
	
\end{document}